\section{Function spaces}\label{functionspaces}
%-----------------------------------------------------------------------------------------------------------------------

In this section we fix notation for function spaces and recall well-known facts.

%-----------------------------------------------------------------------------------------------------------------------
\subsection{H\"older spaces}
%-----------------------------------------------------------------------------------------------------------------------

Let $\Om \subseteq \R^n$ be open and bounded. We~denote by $C^0(\Om)$ the space of continuous complex valued functions on $\Om$.
For $k \in \N \cup \{\infty\}$ (and multi-indices $\ga$) we set
\begin{gather*}
 C^k(\Om)= \big\{f \in \C^\Om \colon \p^\ga f \in C^0(\Om), 0 \le |\ga| \le k\big\},
 \\
 C^k(\overline \Om) = \big\{f \in C^k(\Om) \colon \p^\ga f \text{ has a continuous extension to } \overline \Om, 0 \le |\ga| \le k\big\}.
\end{gather*}
For $\al \in (0,1]$ a function $f \colon \Om \to \C$ belongs to $C^{0,\al}(\overline \Om)$ if it is \emph{$\al$-H\"older continuous}
in $\Om$, i.e.,
\begin{gather*}
\Hoeld_{\al,\Om}(f) := \sup_{x,y \in \Om,\, x \ne y} \frac{|f(x)-f(y)|}{|x-y|^\al} < \infty.
\end{gather*}
If $f$ is \emph{Lipschitz}, i.e., $f \in C^{0,1}\big(\ol \Om\big)$, we write
$\Lip_\Om (f) :=\Hoeld_{1,\Om}(f)$. We~define
\begin{gather*}
 C^{k,\al}\big(\overline \Om\big) =\big \{f \in C^k\big(\overline \Om\big) \colon \p^\ga f \in C^{0,\al}\big(\overline \Om\big), |\ga|\le k\big\},
\end{gather*}
which is a Banach space when provided with the norm
\begin{gather*}
\|f\|_{C^{k,\al}(\overline \Om)}
:= \max_{|\ga| \le k} \sup_{x \in \Om} \big|\p^\ga f(x)\big| + \max_{|\ga|=k} \Hoeld_{\al,\Om}\big(\p^\ga f\big).
\end{gather*}

\subsection{Lebesgue spaces and weak Lebesgue spaces} %\label{Lebesgue}

Let $\Om \subseteq \R^n$ be open and $1 \le p \le \infty$.
Then $L^p(\Om)$ is the Lebesgue space with respect to the $n$-dimensional Lebesgue measure $\cL^n$.
For Lebesgue measurable sets $E \subseteq \R^n$ we denote by
\begin{gather*}
 |E| = \cL^n(E)
\end{gather*}
the $n$-dimensional Lebesgue measure of $E$.
Let $p' := p/(p-1)$ denote the conjugate exponent of $p$
with the convention $1' := \infty$ and $\infty' :=1$.

Let $1 \le p < \infty$ and let us assume that $\Om$ is bounded.
The \emph{weak $L^p$-space} $L_w^p(\Om)$ is the space of all measurable functions $f \colon \Om \to \C$ such that
\begin{gather*}
\|f\|_{p,w,\Om} := \sup_{r> 0} \big( r\, |\{x \in \Om \colon |f(x)| > r\}|^{1/p} \big) < \infty.
\end{gather*}
It will be convenient to \emph{normalize}:
\begin{gather*}
 \|f\|^*_{L^p(\Om)} := |\Om|^{-1/p} \|f\|_{L^p(\Om)},\\
 \|f\|^*_{p,w,\Om} := |\Om|^{-1/p} \|f\|_{p,w,\Om}.
\end{gather*}
Note that $\|1\|^*_{L^p(\Om)} = \|1\|^*_{p,w,\Om} =1$.
For $1 \le q < p < \infty$ we have (cf.\ \cite[Exercise\ 1.1.11]{Grafakos08})
\begin{gather}
 \|f\|^*_{L^q(\Om)} \le \|f\|^*_{L^p(\Om)}, \notag \\
 \|f\|^*_{q,w,\Om} \le \|f\|^*_{L^q(\Om)} \le \bigg(\frac{p}{p-q}\bigg)^{1/q}
 \|f\|^*_{p,w,\Om} \label{eq:qp}
\end{gather}
and hence
$L^p(\Om) \subseteq L_w^p(\Om) \subseteq L^q(\Om) \subseteq L_w^q(\Om)$
with strict inclusions.





We remark that $\|\cdot\|_{p,w,\Om}$ is only a quasinorm: the triangle inequality fails, but for
$f_j \in L_w^p(\Om)$
we still have
\begin{gather*}
\bigg\|\sum_{j=1}^m f_j \bigg\|_{p,w,\Om} \le m \sum_{j=1}^m \|f_j\|_{p,w,\Om}.
\end{gather*}
There exists a norm equivalent to $\|\cdot\|_{p,w,\Om}$ which makes $L_w^p(\Om)$ into a Banach space if $p>1$.

The $L^p_w$-quasinorm is $\si$-subadditive: if
$\Om = \bigcup \Om_j$ is a countable open cover, then
\begin{gather*}
 \|f\|^p_{p,w,\Om} \le \sum_j \|f\|^p_{p,w,\Om_j} \qquad \text{for every}\quad f \in L^p_w(\Om).
\end{gather*}
But it is not $\si$-additive.

\subsection{Sobolev spaces}
For $k \in \N$ and $1 \le p \le \infty$ we consider the Sobolev space
\begin{gather*}
 W^{k,p}(\Om) = \big\{f \in L^p(\Om) \colon \p^\al f \in L^p(\Om), \, 0 \le |\al| \le k\big\},
\end{gather*}
where $\p^\al f$ denote distributional derivatives, with the norm
\begin{gather*}
 \|f\|_{W^{k,p}(\Om)} := \sum_{|\al| \le k} \|\p^\al f\|_{L^p(\Om)}.
\end{gather*}
On bounded intervals $I \subseteq \R$ the Sobolev space $W^{1,1}(I)$
coincides with the space $AC(I)$ of~abso\-lutely continuous functions on~$I$
if we identify each $W^{1,1}$-function with its unique continuous representative.
Recall that a function $f \colon \Om \to \C$ on an open subset $\Om \subseteq \R$
is absolutely continuous ($AC$) if for every $\ep>0$ there exists $\de>0$ such that
for every finite collection of~non-overlapping intervals $(a_i,b_i)$, $i =1,\dots,n$, with
$[a_i,b_i] \subseteq \Om$ we have
\begin{gather*}
 \sum_{i=1}^n |a_i -b_i| < \de \implies \sum_{i=1}^n |f(a_i) -f(b_i)| < \ep.
\end{gather*}
Notice that $W^{1,\infty}(\Om) \cong C^{0,1}\big(\ol \Om\big)$ on Lipschitz domains (or more generally quasiconvex domains)~$\Om$.

We shall also use $W^{k,p}_{\on{loc}}$, $AC_{\on{loc}}$, etc.\ with the obvious meaning.

\subsection{Vector valued functions}

For our problem we need to consider mappings of Sobolev regularity with values in a finite-dimensional complex vector space $V$.
Let us fix a basis $v_1, \dots, v_n$ of $V$ and hence a linear isomorphism $\vh \colon V \to \C^n$. We~say that a mapping $f \colon \Om \to V$ is of Sobolev class $W^{k,p}$ if $\vh \o f$ is of class $W^{k,p}$.
The space $W^{k,p}(\Om, V)$ of all such mappings does not depend on the choice of the basis of $V$.

For $f= (f_1,\dots,f_n) \colon \Om \to \C^n$ we set
\begin{gather} \label{vectorvalued}
 \|f\|_{W^{k,p}(\Om,\C^n)} := \sum_{j=1}^n \|f_j\|_{W^{k,p}(\Om)}.
\end{gather}
If $f \in W^{k,p}(\Om, V)$, $f \ne 0$, and $\vh,\ps \colon V \to \C^n$ are two different basis isomorphisms, then
\begin{gather*}
 c \le \frac{\|\vh \o f\|_{W^{k,p}(\Om,\C^n)}}{\|\ps \o f\|_{W^{k,p}(\Om,\C^n)}} \le C
\end{gather*}
for positive constants $c,C>0$ which depend only on the linear isomorphism $\vh \o \ps^{-1}$. We~will denote by ${\|f\|_{W^{k,p}(\Om,V)}}$ or simply ${\|f\|_{W^{k,p}(\Om)}}$ any of the equivalent norms
$\|\vh \o f\|_{W^{k,p}(\Om,\C^n)}$.


Now suppose that we have a representation $\rh \colon G \to \on{GL}(V)$ of a finite group $G$ on
$V$. By~fixing a Hermitian inner product on $V$ and averaging it over $G$ we obtain a Hermitian inner product
with respect to which the action of~$G$ is unitary. We~could equivalently define
\begin{gather*}
 \|f\|_{W^{k,p}(\Om)} = \|f\|_{W^{k,p}(\Om,V)} := \sum_{|\al| \le k} \bigg(\int_{\Om} \|\p^\al f\|^p \, {\rm d}x\bigg)^{1/p},
 \end{gather*}
 where $\|\cdot\|$ is the norm associated with the $G$-invariant Hermitian inner product. In~that case $\|f\|_{W^{k,p}(\Om,V)}$ is $G$-invariant.


%-----------------------------------------------------------------------------------------------------------------------
\subsection{Extension lemma}
%-----------------------------------------------------------------------------------------------------------------------


The following extension lemma simply follows from the $\C$-valued version proved in \cite{ParusinskiRainer15}.
Similar versions can be found in \cite[Lemma 2.1]{ParusinskiRainerAC}
and \cite[Lemma 3.2]{GhisiGobbino13}.

\begin{Lemma} \label{lem:extend}
 Let $V$ be a finite-dimensional vector space.
 Let $\Om \subseteq \R$ be open and bounded, let $f \colon \Om \to V$ be continuous, $p \ge 1$,
 and set $\Om_0 := \{t \in \Om \colon f(t) \ne 0\}$.
 Assume that $f|_{\Om_0} \in AC_{\on{loc}}(\Om_0,V)$ and $f|_{\Om_0}' \in L^p(\Om_0,V)$.
 Then the distributional derivative of $f$ in $\Om$ is a measurable function $f' \in L^p(\Om,V)$ and
 \begin{gather*} %\label{eq:extend}
 \|f'\|_{L^p(\Om,V)} = \|f|_{\Om_0}'\|_{L^p(\Om_0,V)},
 \end{gather*}
 where the $L^p$-norms are computed with respect to a fixed basis isomorphism.
\end{Lemma}

\section[Finite rotation groups in C]{Finite rotation groups in $\C$}\label{radicals}

Let $C_d \cong \Z/d\Z$ denote the cyclic group of order $d$ and consider its standard action on $\C$ by~rota\-tion.
Then $\C[\C]^{C_d}$ is generated by $\si(z) = z^d$. A lift over $\si$ of a function $f \colon \Om \to \C$ is a~solution of the
equation $z^d = f$.

The solution of the lifting problem in this simple example is completely understood. We~shall see that the general
solution is based on this prototypical case. Interestingly, it is also the case with the worst loss of regularity.

The following theorem is a consequence of a result of Ghisi and Gobbino \cite{GhisiGobbino13}.

\begin{Theorem} \label{cor:radicals}
 Let $d$ be a positive integer and
 let $I \subseteq \R$ be an open bounded interval. Assume that $f \colon I \to \C$ is a continuous function such that
 $f^d = g \in C^{d-1,1}\big(\oI\big)$.
 Then we have $f' \in L^{d'}_w(I)$ and
 \begin{gather} \label{est}
 \|f'\|_{d',w,I} \le
 C(d) \max\Big\{\big(\Lip_{I}\big(g^{(d-1)}\big)\big)^{1/d}|I|^{1/d'},
 \|g'\|_{L^\infty(I)}^{1/d}\Big\}.
 \end{gather}
\end{Theorem}

In other words
any continuous lift $f$ over $\si(z) = z^d$ of a curve in $C^{d-1,1}\big(\ol I,\si(\C)\big) = C^{d-1,1}\big(\ol I\big)$
is absolutely continuous and $f' \in L^{d'}_w(I)$ with the uniform bound~\eqref{est}.

\begin{Remark} \label{optimality}
This result is optimal: in general, $f'$ is not in $L^{d'}$ even if $g$ is real analytic (consider~$g(t) = t$).
On the other hand, if $g$ is only of class $C^{d-1,\be}\big(\overline I\big)$ for every $\be<1$, then $f$
does in~general not need to have bounded variation
in $I$ (see \cite[Example 4.4]{GhisiGobbino13}).
\end{Remark}
